\section{Conclusion}

This project shows that temporal representation is the most important 
factor for accurate bike rental demand forecasting at station level. 
Adding station-specific lag features improved the test RMSE for every 
evaluated model, which confirms that recent demand history carries 
information that weather and calendar features alone cannot provide.

The best final model is Gradient Boosting with lag features, reaching 
a test RMSE of 22.78, MAE of 16.92, and $R^2$ of 0.663. The result is 
not only about which algorithm wins, but about why: the final dataset 
gives the model station identity, weather context, calendar structure, 
and recent demand memory. Other key decisions -- removing leakage-prone 
count columns, reducing redundant weather variables, and using a strict 
chronological split -- ensured that the evaluation reflects real 
forecasting conditions.

Future work should move from daily to hourly forecasting, add event 
and holiday information, and evaluate model performance across all 
stations rather than only the 20 most popular ones. These steps would 
make the model more suitable for real operational planning.